% GENERATED FILE. Edit canonical lessons, then run npm run generate:books.
% canonical-source-hash: fnv1a64:a32cf585730df2ea
% canonical-lessons: MR-C520-nal, MR-C520-galne, MR-C520-galti, MR-C520-gatar, MR-C520-tumbne

\chapter{Pipes and Leaks}
\label{ch:mr-pipes-and-leaks}
\hlchaptermodality{520}

\begin{chapteropening}
\textbf{By the end of this chapter:} \emph{I can talk about water pipes, leaks and blocked drains.}

Complete the last lesson of chapter 520: I can talk about water pipes, leaks and blocked drains.
\end{chapteropening}

\section[\texorpdfstring{\mr{नळ} (naḷ)}{नळ (naḷ)}]{\mr{नळ} (naḷ) — a faucet; a water pipe}

\label{lesson:MR-C520-nal}



\begin{quote}
Before the new one: say the Marathi for darkness, then the Marathi for to light up.
\end{quote}

\subsection*{You'll want to know: \mr{नळ}}
\textbf{\mr{नळ}} — \emph{naḷ} — \textquotedblleft{}a faucet; a water pipe\textquotedblright{}.

\begin{quote}
\mr{आज} \mr{सकाळपासून} \mr{नळाला} \mr{पाणी} \mr{नाही}.
\end{quote}

\begin{quote}
\emph{āj sakāḷpāsūn naḷālā pāṇī nāhī.}
\end{quote}

\begin{quote}
There has been no water in the tap since this morning.
\end{quote}

Masculine.

The tap takes -\mr{ला}: \mr{नळाला} \mr{पाणी} \mr{येतं}, water comes in the tap.

\begin{cousinweb}
\textbf{\mr{नळ}} (\emph{naḷ}) is a tadbhava: CDIAL lists Marathi \emph{naḷ}, \textquotedblleft{}lotus stalk; pipe, conduit\textquotedblright{}, under Sanskrit \emph{naḍa}, \textquotedblleft{}a reed\textquotedblright{} (CDIAL 6936): a hollow stalk and a pipe share one shape.
\end{cousinweb}

\begin{grammarlens}[title={What you've built}]
One more word for describing what you lived and felt.
\end{grammarlens}

\subsection*{Your turn}
\begin{itemize}
\raggedright
  \item \emph{Say it:} \emph{naḷ}
  \item \emph{Say it:} \emph{naḷ}, once more
  \item \emph{Say it:} \emph{āj sakāḷpāsūn naḷālā pāṇī nāhī.} — There has been no water in the tap since this morning.
  \item \emph{Recall:} say the Marathi for darkness, then the Marathi for to light up, then say \emph{naḷ} again
\end{itemize}

\begin{tcolorbox}[breakable,colback=teal!4,colframe=teal!35!black,title={Before you move on}]
What does \mr{नळ} mean? (\textquotedblleft{}A faucet; a water pipe\textquotedblright{}.) Say it once more.
\end{tcolorbox}
\section[\texorpdfstring{\mr{गळणे} (gaḷṇe)}{गळणे (gaḷṇe)}]{\mr{गळणे} (gaḷṇe) — to drip; to leak}

\label{lesson:MR-C520-galne}



\begin{quote}
Before the new one: say the Marathi for to light up, then the Marathi for a faucet.
\end{quote}

\subsection*{You'll want to know: \mr{गळणे}}
\textbf{\mr{गळणे}} — \emph{gaḷṇe} — \textquotedblleft{}to drip; to leak\textquotedblright{}.

\begin{quote}
\mr{तो} \mr{नळ} \mr{रात्रभर} \mr{गळत} \mr{होता}.
\end{quote}

\begin{quote}
\emph{to naḷ rātrabhar gaḷat hotā.}
\end{quote}

\begin{quote}
That tap was dripping all night.
\end{quote}

Past (he …): \mr{तो} \mr{नळ} \mr{गळला}, \emph{to naḷ gaḷlā}.

\begin{cousinweb}
\textbf{\mr{गळणे}} (\emph{gaḷṇe}) is a tadbhava, inherited from Sanskrit \emph{galati}, \textquotedblleft{}drips, trickles\textquotedblright{} (CDIAL 4074).
\end{cousinweb}

\begin{grammarlens}[title={What you've built}]
One more word for describing what you lived and felt.
\end{grammarlens}

\subsection*{Your turn}
\begin{itemize}
\raggedright
  \item \emph{Say it:} \emph{gaḷṇe}
  \item \emph{Say it:} \emph{gaḷṇe}, once more
  \item \emph{Say it:} \emph{to naḷ rātrabhar gaḷat hotā.} — That tap was dripping all night.
  \item \emph{Recall:} say the Marathi for to light up, then the Marathi for a faucet, then say \emph{gaḷṇe} again
\end{itemize}

\begin{tcolorbox}[breakable,colback=teal!4,colframe=teal!35!black,title={Before you move on}]
What does \mr{गळणे} mean? (\textquotedblleft{}To drip; to leak\textquotedblright{}.) Say it once more.
\end{tcolorbox}
\section[\texorpdfstring{\mr{गळती} (gaḷtī)}{गळती (gaḷtī)}]{\mr{गळती} (gaḷtī) — a leak}

\label{lesson:MR-C520-galti}



\begin{quote}
Before the new one: say the Marathi for a faucet, then the Marathi for to drip.
\end{quote}

\subsection*{You'll want to know: \mr{गळती}}
\textbf{\mr{गळती}} — \emph{gaḷtī} — \textquotedblleft{}a leak\textquotedblright{}.

\begin{quote}
\mr{छपरातून} \mr{पाण्याची} \mr{गळती} \mr{होते}.
\end{quote}

\begin{quote}
\emph{chaprātūn pāṇyācī gaḷtī hote.}
\end{quote}

\begin{quote}
Water leaks through the roof.
\end{quote}

Feminine.

Built on \mr{गळणे}: \mr{गळती} \mr{होणे}, there is a leak.

\begin{cousinweb}
\textbf{\mr{गळती}} (\emph{gaḷtī}) is built on \mr{गळणे}, \textquotedblleft{}to drip\textquotedblright{}, inherited from Sanskrit \emph{galati}, \textquotedblleft{}drips\textquotedblright{} (CDIAL 4074).
\end{cousinweb}

\begin{grammarlens}[title={What you've built}]
One more word for describing what you lived and felt.
\end{grammarlens}

\subsection*{Your turn}
\begin{itemize}
\raggedright
  \item \emph{Say it:} \emph{gaḷtī}
  \item \emph{Say it:} \emph{gaḷtī}, once more
  \item \emph{Say it:} \emph{chaprātūn pāṇyācī gaḷtī hote.} — Water leaks through the roof.
  \item \emph{Recall:} say the Marathi for a faucet, then the Marathi for to drip, then say \emph{gaḷtī} again
\end{itemize}

\begin{tcolorbox}[breakable,colback=teal!4,colframe=teal!35!black,title={Before you move on}]
What does \mr{गळती} mean? (\textquotedblleft{}A leak\textquotedblright{}.) Say it once more.
\end{tcolorbox}
\section[\texorpdfstring{\mr{गटार} (gaṭār)}{गटार (gaṭār)}]{\mr{गटार} (gaṭār) — a drain; a gutter}

\label{lesson:MR-C520-gatar}



\begin{quote}
Before the new one: say the Marathi for to drip, then the Marathi for a leak.
\end{quote}

\subsection*{You'll want to know: \mr{गटार}}
\textbf{\mr{गटार}} — \emph{gaṭār} — \textquotedblleft{}a drain; a gutter\textquotedblright{}.

\begin{quote}
\mr{पावसाळ्यात} \mr{रस्त्यावरचं} \mr{गटार} \mr{भरून} \mr{वाहतं}.
\end{quote}

\begin{quote}
\emph{pāvsāḷyāt rastyāvarcaṁ gaṭār bharūn vāhtaṁ.}
\end{quote}

\begin{quote}
In the rainy season the drain on the road overflows.
\end{quote}

Neuter.

\begin{cousinweb}
\textbf{\mr{गटार}} (\emph{gaṭār}) is the word for an open drain; its origin is not settled.
\end{cousinweb}

\begin{grammarlens}[title={What you've built}]
One more word for describing what you lived and felt.
\end{grammarlens}

\subsection*{Your turn}
\begin{itemize}
\raggedright
  \item \emph{Say it:} \emph{gaṭār}
  \item \emph{Say it:} \emph{gaṭār}, once more
  \item \emph{Say it:} \emph{pāvsāḷyāt rastyāvarcaṁ gaṭār bharūn vāhtaṁ.} — In the rainy season the drain on the road overflows.
  \item \emph{Recall:} say the Marathi for to drip, then the Marathi for a leak, then say \emph{gaṭār} again
\end{itemize}

\begin{tcolorbox}[breakable,colback=teal!4,colframe=teal!35!black,title={Before you move on}]
What does \mr{गटार} mean? (\textquotedblleft{}A drain; a gutter\textquotedblright{}.) Say it once more.
\end{tcolorbox}
\section[\texorpdfstring{\mr{तुंबणे} (tumbṇe)}{तुंबणे (tumbṇe)}]{\mr{तुंबणे} (tumbṇe) — to get blocked; to back up}

\label{lesson:MR-C520-tumbne}



\begin{quote}
Before the new one: say the Marathi for a leak, then the Marathi for a drain.
\end{quote}

\subsection*{You'll want to know: \mr{तुंबणे}}
\textbf{\mr{तुंबणे}} — \emph{tumbṇe} — \textquotedblleft{}to get blocked; to back up\textquotedblright{}.

\begin{quote}
\mr{खूप} \mr{पाऊस} \mr{पडला} \mr{आणि} \mr{गटार} \mr{तुंबलं}.
\end{quote}

\begin{quote}
\emph{khūp pāūs paḍlā āṇi gaṭār tumblaṁ.}
\end{quote}

\begin{quote}
It rained heavily and the drain got blocked.
\end{quote}

Past (he …): \mr{तो} \mr{नळ} \mr{तुंबला}, \emph{to naḷ tumblā}.

\begin{cousinweb}
\textbf{\mr{तुंबणे}} (\emph{tumbṇe}) is a tadbhava: CDIAL lists Marathi \emph{tũbṇẽ}, \textquotedblleft{}to accumulate through being obstructed\textquotedblright{}, under Sanskrit \emph{tumbati} (CDIAL 5870).
\end{cousinweb}

\begin{grammarlens}[title={What you've built}]
One more word for describing what you lived and felt.
\end{grammarlens}

\subsection*{Your turn}
\begin{itemize}
\raggedright
  \item \emph{Say it:} \emph{tumbṇe}
  \item \emph{Say it:} \emph{tumbṇe}, once more
  \item \emph{Say it:} \emph{khūp pāūs paḍlā āṇi gaṭār tumblaṁ.} — It rained heavily and the drain got blocked.
  \item \emph{Recall:} say the Marathi for a leak, then the Marathi for a drain, then say \emph{tumbṇe} again
\end{itemize}

\begin{tcolorbox}[breakable,colback=teal!4,colframe=teal!35!black,title={Before you move on}]
What does \mr{तुंबणे} mean? (\textquotedblleft{}To get blocked; to back up\textquotedblright{}.) Say it once more.
\end{tcolorbox}
